\subsection{Direct morphisms}

DEFINITION 1. --- Let E be a topological vector space. A vector subspace of E is said to be direct if it admits a topological complement.

For a vector subspace F of E to be direct, it is necessary and sufficient that there exist a projector of image F in \(\mathcal{L}(\mathrm{E})\) . If \(p\) is such a projector, E is the topological direct sum of F and the kernel of \(p\) . The kernel of \(1_{\mathrm{E}} - p\) is F; if E is separated, then F is therefore closed.

Let F be a direct vector subspace of E and \(E_{1}\) a vector subspace of E containing F; then F is a direct vector subspace of \(E_{1}\) .

Every closed vector subspace of finite codimension in E is a direct summand (EVT, I, p. 15, prop. 3), and every vector subspace F of E contained in the closure of \{0\} is direct (indeed, there exists a projector with image F, and such a projector is necessarily continuous).

PROPOSITION 1. --- Let E be a locally convex space. Every finite-dimensional vector subspace of E is complemented.

Let F be a finite-dimensional vector subspace of E. Let S denote a topological complement of the closure of \(\{0\}\) in F, which is a direct subspace of E. The space \(F_{2} = F \cap S\) is separated and of finite dimension; there therefore exists an integer \(n \geqslant 0\) such that \(F_{2}\) is isomorphic to \(K^{n}\) (EVT, I, p. 14, th. 2). There then exists a continuous linear map \(p: S \to F_{2}\) extending the identity map of \(F_{2}\) . The kernel of p is a topological supplement of \(F_{2}\) in S, and also of F in E.

DEFINITION 2. --- Let E and F be locally convex spaces, and let u be a continuous linear map from E into F. One says that u is a direct morphism if u is a strict morphism whose kernel is a direct vector subspace of E and whose image is a direct vector subspace of F.

Let \(u \in \mathcal{L}(\mathrm{E};\mathrm{F})\) . For the kernel and image of \(u\) to be direct vector subspaces of E and F, respectively, it is necessary and sufficient that there exist decompositions \(\mathrm{E} = \mathrm{E}_1 \oplus \mathrm{E}_2\) and \(\mathrm{F} = \mathrm{F}_1 \oplus \mathrm{F}_2\) into topological direct sums such that \(u\) be represented by a matrix \(\begin{pmatrix} u_1 & 0 \\ 0 & 0 \end{pmatrix}\) where \(u_1 \in \mathcal{L}(\mathrm{E}_1; \mathrm{F}_1)\) is bijective. Since the kernel of \(u\) is then \(\mathrm{E}_2\), the canonical linear map from \(\mathrm{E}_1\) onto \(\mathrm{E}/\mathrm{E}_2\) is an isomorphism, and the image of \(u\) is \(\mathrm{F}_1\), one sees that \(u\) is a strict morphism if and only if \(u_1\) is an isomorphism of \(\mathrm{E}_1\) onto \(\mathrm{F}_1\). Denote by \(v\) the element of \(\mathcal{L}(\mathrm{F};\mathrm{E})\) represented by the matrix \(\begin{pmatrix} u_1^{-1} & 0 \\ 0 & 0 \end{pmatrix}\). Then the linear map \(u \circ v\) is the projector in F with kernel \(\mathrm{F}_2\) and image \(\mathrm{F}_1\), the linear map \(v \circ u\) is the projector in E with kernel \(\mathrm{E}_2\) and image \(\mathrm{E}_1\), and one has \(u \circ v \circ u = u\).

Conversely, one has the following result:

PROPOSITION 2. --- Let E and F be locally convex spaces, let \(u \in \mathcal{L}(\mathrm{E};\mathrm{F})\) and \(v \in \mathcal{L}(\mathrm{F};\mathrm{E})\) . Suppose that \(u = u \circ v \circ u\) . Then \(u\) is a direct morphism. Moreover, \(v \circ u\) is a continuous projector in E with kernel \(\operatorname{Ker}(u)\) and \(u \circ v\) is a continuous projector in F with image \(\operatorname{Im}(u)\) .

Set \(p = v \circ u\) and \(q = u \circ v\) . We have

\[
p ^ {2} = v \circ (u \circ v \circ u) = v \circ u = p, q ^ {2} = (u \circ v \circ u) \circ v = u \circ v = q,
\]

hence p and q are projectors. They are continuous. Let \(E_{1}\) and \(E_{2}\) be the image and the kernel of p, and let \(F_{1}\) and \(F_{2}\) be the image and the kernel of q. Since \(u \circ v \circ u = u\) , we have

\[
\operatorname{Ker} (u) \subset \operatorname{Ker} (v \circ u) \subset \operatorname{Ker} (u \circ v \circ u) = \operatorname{Ker} (u),
\]

hence \(\operatorname{Ker}(u) = \operatorname{Ker}(v \circ u) = \operatorname{Ker}(p) = \operatorname{E}_{2}\) . Similarly, one deduces from the inclusions

\[
\operatorname{Im} (u) \supset \operatorname{Im} (u \circ v) \supset \operatorname{Im} (u \circ v \circ u) = \operatorname{Im} (u)
\]

that the image of \(u\) is \(\mathrm{F}_{1}\) . The spaces \(\mathrm{E}_{2}\) and \(\mathrm{F}_{1}\) are direct vector subspaces of E and F respectively.

The vector space E is the topological direct sum of \(E_{1}\) and \(E_{2}\) , and u defines by restriction a bijective linear map \(u_{1}\) of \(E_{1}\) onto \(F_{1}\) . For every \(x \in E_{1}\) , we have \(v(u(x)) = p(x) = x\) , and the map \(u_{1}^{-1}\) of \(F_{1}\) onto \(E_{1}\) coincides with v on \(F_{1}\) . Consequently, \(u_{1}\) is an isomorphism of \(E_{1}\) onto \(F_{1}\) and u is a strict morphism from E into F.

Let E and F be locally convex spaces and \(u \in \mathcal{L}(\mathrm{E};\mathrm{F})\) . For u to be an injective (resp. surjective) direct morphism, it is necessary and sufficient that there exist a \(v\) in \(\mathcal{L}(\mathrm{F};\mathrm{E})\) such that \(v \circ u = 1_{E}\) (resp. \(u \circ v = 1_{F}\) ) (cf. TG, III, p. 47 and 48).

PROPOSITION 3. --- Let E and F be Banach spaces. The following sets are open subsets of the Banach space \(\mathcal{L}(\mathrm{E};\mathrm{F})\) :

\begin{enumerate}
\def\labelenumi{\alph{enumi})}
\item
  The set \(\mathcal{I}(\mathrm{E};\mathrm{F})\) of isomorphisms from E onto F;
\item
   The set \(\mathcal{M}\mathcal{D}(\mathrm{E};\mathrm{F})\) of injective direct morphisms from E into F, and more precisely, for every closed vector subspace \(F_{1}\) of F, the set \(\mathcal{M}_{\mathrm{F}_{1}}(\mathrm{E};\mathrm{F})\) of elements of \(\mathcal{M}\mathcal{D}(\mathrm{E};\mathrm{F})\) whose image is a topological supplement of \(F_{1}\) ;
\item
   The set \(\mathcal{E}\mathcal{D}(\mathrm{E};\mathrm{F})\) of surjective direct morphisms from E into F, and more precisely, for every closed vector subspace \(\mathrm{E}_1\) of E, the set \(\mathcal{E}_{\mathrm{E}_1}(\mathrm{E};\mathrm{F})\) of elements of \(\mathcal{E}\mathcal{D}(\mathrm{E};\mathrm{F})\) whose kernel is a topological supplement of \(\mathrm{E}_1\) .
\end{enumerate}

Moreover, the map \(u \mapsto u^{-1}\) of \(\mathcal{I}(\mathrm{E};\mathrm{F})\) onto \(\mathcal{I}(\mathrm{F};\mathrm{E})\) is analytic.

By definition, \(\mathcal{I}(\mathrm{E};\mathrm{E})\) is the set of invertible elements of the complete normed algebra \(\mathcal{L}(\mathrm{E})\) . By TG, IX, p. 40, prop. 14, it is an open subset of \(\mathcal{L}(\mathrm{E})\) . The map \(u\mapsto u^{-1}\) from \(\mathcal{I}(\mathrm{E};\mathrm{E})\) into \(\mathcal{I}(\mathrm{E};\mathrm{E})\) is analytic (LIE, III, § 1, n° 1, example 2).

If the set \(\mathcal{I}(\mathrm{E};\mathrm{F})\) is empty, it is open. Otherwise, let \(u_0\) be an isomorphism from E onto F. The map \(v\mapsto u_0\circ v\) is then an isomorphism of \(\mathcal{L}(\mathrm{E})\) onto \(\mathcal{L}(\mathrm{E};\mathrm{F})\) which transforms \(\mathcal{I}(\mathrm{E};\mathrm{E})\) into \(\mathcal{I}(\mathrm{E};\mathrm{F})\) . The set \(\mathcal{I}(\mathrm{E};\mathrm{F})\) is therefore open in \(\mathcal{L}(\mathrm{E};\mathrm{F})\) . Moreover, if \(u = u_0\circ v\) is an element of \(\mathcal{I}(\mathrm{E};\mathrm{F})\) , we have \(u^{-1} = v^{-1}\circ u_0^{-1}\) , and the map \(u\mapsto u^{-1}\) of \(\mathcal{I}(\mathrm{E};\mathrm{F})\) onto \(\mathcal{I}(\mathrm{F};\mathrm{E})\) is therefore analytic.

Let \(F_{1}\) be a closed vector subspace of F. For every \(u \in \mathcal{L}(E;F)\) , let \(\overline{u}\) be the element of \(\mathcal{L}(E \times F_{1};F)\) defined by \(\overline{u}(x,y) = u(x) + y\) . The map \(u \mapsto \overline{u}\) from \(\mathcal{L}(E;F)\) into \(\mathcal{L}(E \times F_{1};F)\) is continuous, and \(\mathcal{M}_{\mathrm{F}_{1}}(E;F)\) is the set of elements u of \(\mathcal{L}(E;F)\) such that \(\overline{u}\) belongs to \(\mathcal{I}(E \times F_{1};F)\) . Since \(\mathcal{I}(E \times F_{1};F)\) is open in \(\mathcal{L}(E \times F_{1};F)\)

by what precedes, the set \(\mathcal{M}_{\mathrm{F}_1}(\mathrm{E};\mathrm{F})\) is open in \(\mathcal{L}(\mathrm{E};\mathrm{F})\) . The same is true of \(\mathcal{MD}(\mathrm{E};\mathrm{F})\) , which is the union of the \(\mathcal{M}_{\mathrm{F}_1}(\mathrm{E};\mathrm{F})\) when \(\mathrm{F}_1\) ranges over the set of closed vector subspaces of F.

Let \(E_{1}\) be a closed vector subspace of E, and let p be the canonical map of E onto the quotient Banach space \(E/E_{1}\) . For an element \(u\) of \(\mathcal{L}(E;F)\) to belong to \(\mathcal{E}_{\mathrm{E}_{1}}(\mathrm{E};\mathrm{F})\), it is necessary and sufficient that the map \((u,p)\) from E into \(F\times E/E_{1}\) belong to \(\mathcal{I}(E;F\times E/E_{1})\). As before, we deduce that \(\mathcal{E}_{\mathrm{E}_{1}}(\mathrm{E};\mathrm{F})\) is open in \(\mathcal{L}(\mathrm{E};\mathrm{F})\); the same is true of \(\mathcal{E}\mathcal{D}(\mathrm{E};\mathrm{F})\), which is the union of the \(\mathcal{E}_{\mathrm{E}_{1}}(\mathrm{E};\mathrm{F})\).

